\documentclass[10pt,a4paper]{amsart}
\usepackage{amsmath,amssymb,amsthm,graphicx}
% Retained preamble for the Stacks Project chapter {CH} (shared project preamble).
% The chapter file pinned-original-source.tex only carries \input{{preamble}} on its first line,
% so no preamble text lives inside the chapter itself.  The Stacks Project uses one single shared
% preamble for every chapter: the snapshot stages it three times (sources/stacks-chow/preamble.tex,
% sources/stacks-etale-cohomology/preamble.tex and focused-seeds/AGP/sources/stacks-snapshot-20261011/
% preamble.tex) and those copies are identical in size.  The block below is copied mechanically,
% line for line, from sources/stacks-chow/preamble.tex lines 14-18 and 154-217: the Xy-pic setup and
% the native theorem declarations together with the mathematical macros.  Nothing was added, removed
% or redefined; the wrapper class and the amsmath/amssymb/amsthm/graphicx base packages are supplied
% by the assembler.
% For commutative diagrams we use Xy-pic
% Chapter: stacks-cohomology.tex -- Cohomology of Sheaves.
\usepackage[all]{xy}

% We use 2cell for 2-commutative diagrams.
\xyoption{2cell}

% Theorem environments.
%
\theoremstyle{plain}
\newtheorem{theorem}[subsection]{Theorem}
\newtheorem{proposition}[subsection]{Proposition}
\newtheorem{lemma}[subsection]{Lemma}

\theoremstyle{definition}
\newtheorem{definition}[subsection]{Definition}
\newtheorem{example}[subsection]{Example}
\newtheorem{exercise}[subsection]{Exercise}
\newtheorem{situation}[subsection]{Situation}

\theoremstyle{remark}
\newtheorem{remark}[subsection]{Remark}
\newtheorem{remarks}[subsection]{Remarks}

\numberwithin{equation}{subsection}

% Macros
%
\def\lim{\mathop{\mathrm{lim}}\nolimits}
\def\colim{\mathop{\mathrm{colim}}\nolimits}
\def\Spec{\mathop{\mathrm{Spec}}}
\def\Hom{\mathop{\mathrm{Hom}}\nolimits}
\def\Ext{\mathop{\mathrm{Ext}}\nolimits}
\def\SheafHom{\mathop{\mathcal{H}\!\mathit{om}}\nolimits}
\def\SheafExt{\mathop{\mathcal{E}\!\mathit{xt}}\nolimits}
\def\Sch{\mathit{Sch}}
\def\Mor{\mathop{\mathrm{Mor}}\nolimits}
\def\Ob{\mathop{\mathrm{Ob}}\nolimits}
\def\Sh{\mathop{\mathit{Sh}}\nolimits}
\def\NL{\mathop{N\!L}\nolimits}
\def\CH{\mathop{\mathrm{CH}}\nolimits}
\def\proetale{{pro\text{-}\acute{e}tale}}
\def\etale{{\acute{e}tale}}
\def\QCoh{\mathit{QCoh}}
\def\Ker{\mathop{\mathrm{Ker}}}
\def\Im{\mathop{\mathrm{Im}}}
\def\Coker{\mathop{\mathrm{Coker}}}
\def\Coim{\mathop{\mathrm{Coim}}}

% Boxtimes
%
\DeclareMathSymbol{\boxtimes}{\mathbin}{AMSa}{"02}

%
% Macros for moduli stacks/spaces
%
\def\QCohstack{\mathcal{QC}\!\mathit{oh}}
\def\Cohstack{\mathcal{C}\!\mathit{oh}}
\def\Spacesstack{\mathcal{S}\!\mathit{paces}}
\def\Quotfunctor{\mathrm{Quot}}
\def\Hilbfunctor{\mathrm{Hilb}}
\def\Curvesstack{\mathcal{C}\!\mathit{urves}}
\def\Polarizedstack{\mathcal{P}\!\mathit{olarized}}
\def\Complexesstack{\mathcal{C}\!\mathit{omplexes}}
% \Pic is the operator that assigns to X its picard group, usage \Pic(X)
% \Picardstack_{X/B} denotes the Picard stack of X over B
% \Picardfunctor_{X/B} denotes the Picard functor of X over B
\def\Pic{\mathop{\mathrm{Pic}}\nolimits}
\def\Picardstack{\mathcal{P}\!\mathit{ic}}
\def\Picardfunctor{\mathrm{Pic}}
\def\Deformationcategory{\mathcal{D}\!\mathit{ef}}

% --- appended 2026-10-12 (Chief): the project own definitions of the reference/slogan/history
% display environments, copied mechanically from the shared Stacks preamble lines 8-12
% (\usepackage{verbatim} plus the three \newenvironment lines, which the chapter preamble keeps
% verbatim from the project preamble). Nothing added, removed or redefined: these environments
% carry no mathematical content and the project itself typesets them as comments. Without these
% three lines a counted statement whose own text uses them fails with Environment slogan undefined.
\usepackage{verbatim}
\newenvironment{reference}{\comment}{\endcomment}
\newenvironment{slogan}{\comment}{\endcomment}
\newenvironment{history}{\comment}{\endcomment}
\usepackage{cleveref}
\crefname{originaldefinition}{original definition}{original definitions}
\Crefname{originaldefinition}{Original definition}{Original definitions}
\makeatletter
\newlabel{definition-pseudo-coherent}{{08CB}{source}{}{sourcepointer-1}{}}
\newlabel{definition-pseudo-coherent@cref}{{[originaldefinition][0][]08CB}{[1][0][]source}{}{}{}}
\makeatother
\crefname{originallemma}{original lemma}{original lemmas}
\Crefname{originallemma}{Original lemma}{Original lemmas}
\makeatletter
\newlabel{lemma-pseudo-coherent-independent-representative}{{08CC}{source}{}{sourcepointer-2}{}}
\newlabel{lemma-pseudo-coherent-independent-representative@cref}{{[originallemma][0][]08CC}{[1][0][]source}{}{}{}}
\makeatother
\crefname{originalreference}{original reference}{original references}
\Crefname{originalreference}{Original reference}{Original references}
\makeatletter
\newlabel{derived-lemma-negative-vanishing}{{05TC}{source}{}{sourcepointer-3}{}}
\newlabel{derived-lemma-negative-vanishing@cref}{{[originalreference][0][]05TC}{[1][0][]source}{}{}{}}
\makeatother
\crefname{originallemma}{original lemma}{original lemmas}
\Crefname{originallemma}{Original lemma}{Original lemmas}
\makeatletter
\newlabel{lemma-strictly-perfect-pullback}{{09U6}{source}{}{sourcepointer-4}{}}
\newlabel{lemma-strictly-perfect-pullback@cref}{{[originallemma][0][]09U6}{[1][0][]source}{}{}{}}
\makeatother
\title{Pseudo-coherence is preserved by derived pullback}
\author{The Stacks Project Authors}
\date{stacks-project:cohomology.tex:local-snapshot-20261011; GNU FDL 1.2 or later}
\begin{document}
\maketitle
\noindent\textit{Editorial note.} All delivered mathematical text is the author's own, copied verbatim from the pinned Stacks Project chapter file cohomology.tex; no statement or proof step was rewritten, reordered or removed.\par\medskip
\noindent\textit{Editorial note.} The counted span cites the lemmas recorded in assembly.external\_references with canonical tags as navigation pointers.\par\medskip
\noindent\textit{Allowed standard prerequisites (editorial).}\begin{itemize}
\item Pseudo-coherence of representatives. Cohomology of Sheaves, Lemma cohomology-lemma-pseudo-coherent-independent-representative. Reduce to the pulled-back presentations.
\item Pullback of strictly perfect complexes. Cohomology of Sheaves, Lemma cohomology-lemma-strictly-perfect-pullback. f\_i\textasciicircum{}*E\_i\textasciicircum{}* is strictly perfect.
\item Vanishing for left derived pullback. Derived Categories, Lemma derived-lemma-negative-vanishing. H\textasciicircum{}j(Lf\_i\textasciicircum{}*C) = 0 for j >= m.
\end{itemize}\medskip
\ifdefined\hypertarget\hypertarget{sourcepointer-1}{}\fi
\noindent\textit{Source reference.} Original definition 08CB: sources/stacks-chow/tags.txt, source lines 10596--10596. Cross-chapter standard prerequisite cited by the author inside the counted span; it is used as a pointer and its own proof is not part of this item.\par\medskip
\ifdefined\hypertarget\hypertarget{sourcepointer-2}{}\fi
\noindent\textit{Source reference.} Original lemma 08CC: sources/stacks-chow/tags.txt, source lines 10597--10597. Cross-chapter standard prerequisite cited by the author inside the counted span; it is used as a pointer and its own proof is not part of this item.\par\medskip
\ifdefined\hypertarget\hypertarget{sourcepointer-3}{}\fi
\noindent\textit{Source reference.} Original reference 05TC: sources/stacks-chow/tags.txt, source lines 7476--7476. Cross-chapter standard prerequisite cited by the author inside the counted span; it is used as a pointer and its own proof is not part of this item.\par\medskip
\ifdefined\hypertarget\hypertarget{sourcepointer-4}{}\fi
\noindent\textit{Source reference.} Original lemma 09U6: sources/stacks-chow/tags.txt, source lines 12414--12414. Cross-chapter standard prerequisite cited by the author inside the counted span; it is used as a pointer and its own proof is not part of this item.\par\medskip
\begin{lemma}
\label{lemma-pseudo-coherent-pullback}
Let $f : (X, \mathcal{O}_X) \to (Y, \mathcal{O}_Y)$
be a morphism of ringed spaces. Let $E$ be an object of
$D(\mathcal{O}_Y)$. If $E$ is $m$-pseudo-coherent,
then $Lf^*E$ is $m$-pseudo-coherent.
\end{lemma}

\begin{proof}
Represent $E$ by a complex $\mathcal{E}^\bullet$ of $\mathcal{O}_Y$-modules
and choose an open covering $Y = \bigcup V_i$ and
$\alpha_i : \mathcal{E}_i^\bullet \to \mathcal{E}^\bullet|_{V_i}$
as in Definition \ref{definition-pseudo-coherent}.
Set $U_i = f^{-1}(V_i)$.
By Lemma \ref{lemma-pseudo-coherent-independent-representative}
it suffices to show that $Lf^*\mathcal{E}^\bullet|_{U_i}$ is
$m$-pseudo-coherent.
Choose a distinguished triangle
$$
\mathcal{E}_i^\bullet \to
\mathcal{E}^\bullet|_{V_i} \to
C \to
\mathcal{E}_i^\bullet[1]
$$
The assumption on $\alpha_i$ means exactly that the cohomology sheaves
$H^j(C)$ are zero for all $j \geq m$. Denote $f_i : U_i \to V_i$
the restriction of $f$. Note that $Lf^*\mathcal{E}^\bullet|_{U_i} = 
Lf_i^*(\mathcal{E}|_{V_i})$. Applying $Lf_i^*$ we obtain
the distinguished triangle
$$
Lf_i^*\mathcal{E}_i^\bullet \to
Lf_i^*\mathcal{E}|_{V_i} \to
Lf_i^*C \to
Lf_i^*\mathcal{E}_i^\bullet[1]
$$
By the construction of $Lf_i^*$ as a left derived functor we see that
$H^j(Lf_i^*C) = 0$ for $j \geq m$ (by the dual of Derived Categories, Lemma
\ref{derived-lemma-negative-vanishing}). Hence $H^j(Lf_i^*\alpha_i)$ is an
isomorphism for $j > m$ and $H^m(Lf^*\alpha_i)$ is surjective.
On the other hand,
$Lf_i^*\mathcal{E}_i^\bullet = f_i^*\mathcal{E}_i^\bullet$.
is strictly perfect by Lemma \ref{lemma-strictly-perfect-pullback}.
Thus we conclude.
\end{proof}

\end{document}
